\begin{question}
\noindent $P$, $Q$, $R$ and $S$ are points on the circumference of a circle, forming the cyclic quadrilateral $PQRS$.

\noindent Angle $P = (3x + 15)^\circ$ and angle $R = (2x + 5)^\circ$.

\begin{center}
\begin{tikzpicture}[scale=2]
    \coordinate (P) at (150:1);
    \coordinate (Q) at (60:1);
    \coordinate (R) at (-30:1);
    \coordinate (S) at (220:1);

    \draw[name path=circ] (0,0) circle (1);
    \draw[line join=bevel] (P) -- (Q) -- (R) -- (S) -- cycle;

    \node[above left] at (P) {$P$};
    \node[above right] at (Q) {$Q$};
    \node[below right] at (R) {$R$};
    \node[below left] at (S) {$S$};

    \pic [draw, "$3x+15$", angle eccentricity=1.6, angle radius=0.4cm] {angle = S--P--Q};
    \pic [draw, "$2x+5$", angle eccentricity=1.6, angle radius=0.4cm] {angle = Q--R--S};
\end{tikzpicture}
\end{center}

\noindent Work out the size of angle $R$.

\vspace{4.5cm}

\begin{flushright}
\answerequnit{R}{$^\circ$}{3}
\end{flushright}
\end{question}
